\subsection{Оценка решения}

В данном разделе будет произведена оценка решения.

\subsubsection{Пространственная и временная сложность}

Рассмотрим алгоритмическую сложность \cite{clrs_algorithms} решения. Размеры трассы и количество итераций при проходе по ней пропорциональны количеству инструкций в трассе $N = |V|$.

Рассмотри временную и пространственную сложность стадий обработки трассы:

\begin{itemize}
    \item Генерация карты перемещений

    Временная сложность - $O(N)$ (полный проход по трассе).
    Пространственная сложность - $O(N)$ (агрегация данных с записей трассы).

    \item Сбор ГПД

    Временная сложность - $O(N)$ (полный проход по трассе).
    Пространственная сложность - $O(N + K)$. Где $K = |E|$ - количество ребер ГПД.

    \item Генерация модификаций трассы

    На данной стадии производится два обхода ГПД: построение обратных путей и прямой обход. При этом ни одно ребро не обходится дважды, так как обработанные ребра отмечаются маркерами. Временная сложность обходов графа - $O(N + K)$. Также в конце производится симуляционный проход по модифицированной трассе, временная сложность - $O(N)$.

    Результирующая временная сложность - $O(N + K)$

    На данной стадии генерируются модификации трассы, количество которых пропорционально количеству инструкций и количеству ребер ГПД. Пространственная сложность - $O(N + K)$.

    \item Модификация трассы

    Временная сложность - $O(N)$ (полный проход по трассе).
    Пространственная сложность - $O(1)$.
\end{itemize}

Заметим, что выполняется асимптотическая оценка: $K = O(N)$. Таким образом, пространственная и временная сложность обработки трассы $O(N + K) = O(N)$.

\subsubsection{Оценка репрезентативности трасс}

Для оценки репрезентативности выходных трасс будет использоваться статистика произведенных модификаций. Рассмотрим различные модификации трассы и их влияние на репрезентативность:

\begin{itemize}
    \item Замена адреса для операции с памятью - не влияет на репрезентативность, так как не требует компенсации при дальнейшей обработке выходной трассы.

    \item Замена значения для записи в память - не требует компенсации (поток данных консистентен).

    \item Замена значения для чтения из памяти - может влиять на репрезентативность в некоторых сценариях. При чтении значения, не соответствующего текущему образу памяти рассматриваемого потока, может потребоваться компенсация, аналогичная компенсации для внешней модификации памяти. Заметим, что если в оригинальной трассе уже присутствовала внешняя модификация памяти, модификация не повлияет на репрезентативность.

    \item Замена значения регистра в начальном состоянии - не требует дополнительной компенсации.

    \item Замена значения для записи в регистр общего назначения - не требует дополнительной компенсации.

    \item Добавление записи в регистр общего назначения - требует компенсации, будет снижать репрезентативность.
\end{itemize}

Таким образом, для оценки влияния на репрезентативность трасс будем использовать кумулятивную метрику $N_{comp}$ - количество модификаций, приводящих к появлению дополнительных компенсаций при обработке трассы. Обозначим также $P_{comp} = \frac{N_{comp}}{|V|}$ - количество компенсируемых модификаций на инструкцию трассы.

\subsubsection{Профилировка решения по памяти} \label{sec:memory_profiling}

Для получения грубой оценки потребления памяти различными компонентами использовалась утилита heaptrack \cite{heaptrack_documentation}, позволяющая отследить объемы выделений памяти различными участками кода. Данные heaptrack для набора тестовых трасс показывают, что в коде библиотеки Boost Graph выделяется больше 90 процентов используемой памяти.

\subsubsection{Тестовые запуски на целевых трассах} \label{sec:real_trace_tests}

В данном разделе будут представлены результаты тестовых запусков на целевых трассах заказчика (таблица \ref{tab:real_traces_tests}). Рассматриваются следующие метрики:

\begin{itemize}
    \item $N$ - количество инструкций в трассе.
    \item $N_{comp}$ - количество компенсируемых модификаций.
    \item $P_{comp} = \frac{N_{comp}}{N}$ - количество компенсируемых модификаций на инструкцию трассы.
    \item $RSS_{max}$ - максимальный объем резидентной памяти (resident set size) \cite{gnu_time_memory_resources,linux_proc_filesystem}, занятый при модификации трассы.
    \item $P_{RSS} = \frac{RSS_{max}}{N}$ - объем резидентной памяти на инструкцию трассы.
    \item $T$ - время обработки трассы.
    \item $V = \frac{N}{T}$ - средняя скорость обработки (инструкций в секунду).
\end{itemize}

\begin{table}[htbp]
\centering \small

\caption{Результаты тестовых запусков на целевых трассах заказчика}
\label{tab:real_traces_tests}

\begin{tabular}{|p{3.6cm}|l|l|l|l|l|l|l|} \hline
    \textbf{Трасса}         & $\mathbf{N}$ & $\mathbf{N_{comp}}$ & $\mathbf{P_{comp}}$ & $\mathbf{RSS_{max}},$ & $\mathbf{P_{RSS}},$         & $\mathbf{T},\,\texttt{С}$ & $V,\,\frac{1}{\texttt{С}}$ \\
    \textbf{(тип нагрузки)} &              &                     &                     & $\texttt{Гб}$         & $\frac{\texttt{Гб}}{10^6}$  &                           &                               \\ \hline

    LogicTrace1 \linebreak (вычислительная) & $1.7 \cdot 10^{7}$ & $3.7 \cdot 10^{5}$ & $2.1 \cdot 10^{-2}$ & $7.65$  & $0.44$ & $162.5$ & $1.07 \cdot 10^{5}$ \\ \hline

    LogicTrace2 \linebreak (вычислительная) & $1.6 \cdot 10^{7}$ & $4.0 \cdot 10^{5}$ & $2.5 \cdot 10^{-2}$ & $6.91$  & $0.44$ & $149.3$ & $1.06 \cdot 10^{5}$ \\ \hline

    LogicTrace3 \linebreak (вычислительная) & $4.1 \cdot 10^{7}$ & $9.5 \cdot 10^{4}$ & $2.3 \cdot 10^{-3}$ & $16.71$ & $0.41$ & $408.7$ & $1.00 \cdot 10^{5}$ \\ \hline

    LogicTrace4 \linebreak (вычислительная) & $1.6 \cdot 10^{7}$ & $7.4 \cdot 10^{3}$ & $4.5 \cdot 10^{-4}$ & $6.91$  & $0.43$ & $160.0$ & $1.05 \cdot 10^{5}$ \\ \hline

    LogicTrace5 \linebreak (вычислительная) & $3.4 \cdot 10^{7}$ & $1.3 \cdot 10^{3}$ & $4.1 \cdot 10^{-5}$ & $13.38$ & $0.40$ & $321.3$ & $1.04 \cdot 10^{5}$ \\ \hline

    GfxTrace1 \linebreak (графика)          & $1.0 \cdot 10^{7}$ & $7.0 \cdot 10^{3}$ & $6.7 \cdot 10^{-4}$ & $4.47$  & $0.43$ & $99.3$  & $1.05 \cdot 10^{5}$ \\ \hline
\end{tabular}

\end{table}

Обозначим: $M$ - количество тестовых трасс. Рассчитаем средние значения для $\mathbf{P_{comp}}$, $\mathbf{P_{RSS}}$ и $V$.

Для $\mathbf{P_{comp}}$ приведено среднее арифметическое - доля компенсированных инструкций характеризует качество выходной трассы, подобное усреднение позволит оценить долю компенсаций в случайно выбранной трассе. При этом не учитываются размеры тестовых трасс, так как на долю компенсаций в первую очередь влияют особенности потоков данных и управления. Также для $\mathbf{P_{comp}}$ приведен 50-й процентиль (медиана), так как на ряде тестовых трасс доля компенсаций существенно выше, что может говорить о наличии большого количества еще не поддержанных краевых случаев. \\

$\displaystyle \overline{P_{comp}} = \frac{1}{M} \sum_{i=1}^{M} P_{comp_i} = 8.3 \cdot 10^{-3} $ \\

$ \displaystyle P_{50}\left(P_{comp}\right) = 6.7 \cdot 10^{-4} $ \\

$\mathbf{P_{RSS}}$ также усреднено с помощью среднего арифметического - потребление памяти линейно от количества инструкций (остальные накладные расходы пренебрежимо малы). Значит потребление памяти на одну инструкцию также будет в первую очередь зависеть от свойств потоков управления и данных (например от количества ребер на одну инструкцию), но не от размеров трасс. Подобная оценка позволит оценить максимальное суммарное количество инструкций в наборе трасс, который можно обрабатывать одновременно в условиях ограниченной памяти. \\

$\displaystyle \overline{P_{RSS}} = \frac{1}{M} \sum_{i=1}^{M} P_{RSS_i} = 0,42 \frac{\texttt{Гб}}{10^6} $ \\

Для $V$ приведено взвешенное гармоническое среднее, являющееся средней скоростью обработки всего массива тестовых трасс. \\

$\displaystyle \overline{V} = \frac{\displaystyle \sum_{i=1}^{M} N_i}{\displaystyle \sum_{i=1}^{M} T_i} = 103334 \frac{1}{\texttt{С}}$ \\
